\documentclass[10pt,a4paper]{amsart}
\usepackage[margin=19mm]{geometry}
\usepackage{amsmath,amssymb,mathtools}
\usepackage[scr=rsfs,cal=euler]{mathalfa}
\usepackage{xfrac}
\usepackage[unicode,hidelinks,bookmarks=false]{hyperref}
\newcommand{\ie}{i.e.\ }
\newcommand{\cf}{cf.\ }
\newcommand{\resp}{respectively\ }
\newcommand{\Subsection}{\subsection}
\providecommand{\nobreakdash}{\nobreak\hbox{-}}
\setlength{\emergencystretch}{2em}
\multlinegap0pt
\usepackage{bm}
\usepackage{bbm}
\usepackage{dsfont}
\usepackage{xspace}
\usepackage{cancel}
\usepackage{mathtools}
\usepackage{amsthm}
\makeatletter
\@ifundefined{thm}{\newtheorem{thm}{Thm}}{}
\@ifundefined{lem}{\newtheorem{lem}[thm]{Lemma}}{}
\@ifundefined{theorem}{\newtheorem{theorem}{Theorem}}{}
\makeatother
\title[On the Area of Immersed Minimal Annuli in a Slab]{On the Area of Immersed Minimal Annuli in a Slab}
\author{Elham Matinpour}
\date{Source: arXiv:2410.15535v2 (CC BY 4.0)}
\begin{document}
\maketitle

\noindent\emph{Source and licence.} On the Area of Immersed Minimal Annuli in a Slab. Version arXiv:2410.15535v2 (CC BY 4.0), distributed under the Creative Commons Attribution 4.0 International licence, as recorded in the arXiv licence metadata for this exact version and shown on the abstract page https://arxiv.org/abs/2410.15535v2. Full rights notice: licenses/arxiv-2410.15535-CC-BY-4.0.txt.\par\smallskip

\noindent\textit{Editorial scope.} Counted unit: the Theorem whose source label is \texttt{thm: existence for L''<4L} in the article (source0.tex lines 572--577, source label \texttt{thm: existence for L''<4L}), delivered together with the article's own complete original proof of it (source0.tex lines 578--673, one \texttt{proof} environment of 96 lines). No mathematics is written, repaired or re-derived: statement and proof are the article's own text line for line. Required context retained uncounted: lem: A* implies int G =0 (lines 492--497). Every definition, display and label of those lines is retained unchanged. Omitted from this package and not redistributed: 1--491, 498--571, 674--1150. Nothing inside a retained excerpt is omitted or altered: every retained body is delivered exactly as the article's lines are written, so no omission receipt of any kind accompanies this package. Citations: the retained excerpts carry no citation command at all, so no bibliography entry of the article is transcribed and no reference is redistributed. Reference adaptation: none was needed and none was made; every reference inside the delivered unit resolves inside it, so no unresolved reference or citation is produced. Printed slips kept literal: no printed word, symbol, hypothesis or number of the counted statement or of its proof was altered. Layout and class: the article's own class is \texttt{amsart} with packages including epsfig, color, blindtext, graphicx, enumitem, url, amssymb, import, comment, esint, xypic, geometry, hyperref; the wrapper is the lane's pinned \texttt{amsart} editorial wrapper at 10pt on A4, which restates the theorem-like environments the retained text needs and adds the article's own macro definitions that the retained text needs (none), with extra wrapper packages (bm, bbm, dsfont, xspace, cancel, mathtools). The delivered numbering is the unit's own. Assets: no retained span contains a figure environment, a graphics-inclusion command, a PDF-page inclusion, a table or a TikZ picture; no image file is copied into this package, the package declares no asset dependency and no image file is redistributed with it. Licence: version arXiv:2410.15535v2 of this article is distributed under the Creative Commons Attribution 4.0 International licence, as recorded in the arXiv licence metadata for this exact version and shown on the abstract page https://arxiv.org/abs/2410.15535v2. A full rights notice is delivered at licenses/arxiv-2410.15535-CC-BY-4.0.txt. Characters: every retained excerpt is pure ASCII, so the delivered engine, XeLaTeX with its default Latin Modern text font, reports no missing character.
\begin{lem} \label{lem: A* implies int G =0}
    If $\bold{X} \in \mathcal{A}_V^*(\Omega ,2k)$ for $k=1,2,3,\cdots$, then 
    $$
    \int_0^{2\pi} G_-(z)d\theta =0,\ \ \text{and}\  \int_0^{2\pi} G_+(z)d\theta =0.
    $$ 
\end{lem}

\begin{theorem} \label{thm: existence for L''<4L}
    There exist minimal surfaces in $\mathcal{A}_V(\Omega,2)\backslash \mathcal{A}^*(\Omega,2)$ for which
    $$
    L''(t)<4L(t),\ \ \forall t
    $$
\end{theorem}

\begin{proof}
    We identify examples of such surfaces by introducing perturbations to the extended Gauss map of a 2-fold cover of a catenoid.\\
 Define the slab $\Omega$ as outlined in the introduction, where $\Omega$ represents the open space between two parallel planes, denoted as $P_{\pm} \subset \mathbb{R}^3$, with $P_{\pm} =\{ x_3 =h_{\pm} \}$. Here, $h_-=-h_+$ with $h_+$ being a positive number. The surface $\bold{F} \in \mathcal{A}_V(\Omega ,2)$ is parameterized by the proper conformal, harmonic immersion $\bold{F}:A_{R_0,R_1} \rightarrow \mathbb{R}^3$, where $A_{R_0,R_1}=\{ z\in \mathbb{C} ;\  R_0< \vert z\vert <R_1 \}$ is an open annular domain. Moreover, $\bold{F}(A_{R_0,R_1}) \subset \Omega$ and $\bar{\Sigma} \backslash \Sigma \subset \partial \Omega$, where $\partial \Omega =  P_+ \cup P_- $. \\
Let $C_2$ be the 2-fold cover of a waist of catenoid parameterized by $\bold{F} \in \mathcal{A}_V(\Omega , 2)$, and assume the following Weierstrass data corresponding to $C_2$:
     $$
     g(z)=\frac{d_1}{z^2} \ \text{and}\ f(z)=d_2z,
     $$
     the corresponding representation is
     $$
     \begin{cases}
         \phi_1 = \frac{d_2z}{2} (1-\frac{d_1^2}{z^4})=\frac{d_2}{2z^3}(z^4-d_1^2) \\
         \phi_2 = \frac{id_2z}{2} (1+\frac{d_1^2}{z^4})=\frac{id_2}{2z^3}(z^4+d_1^2)  \\
         \phi_3 = \frac{d_1d_2}{z}
     \end{cases}
     $$
moreover, $F_-(z)=\frac{\psi_3}{g}=c_1z^2$ and $F_+(z)=\psi_3 g=c_2z^{-2}$.
Consider a perturbation of the Gauss map of $C_2$ given by
    $$
    g(z)=\frac{c_2 z^{-1}+\epsilon_2 +\delta_2 z}{c_1z+\epsilon_1 + \delta_1 z^{-1}},
    $$
   where $\epsilon_1, \delta_1, \epsilon_2, \delta_2$ are sufficiently small numbers, and they satisfy $\epsilon_1^2+2\delta_1c_1=0$ and $\epsilon_2^2+2\delta_2 c_2=0$. 
   By choosing these numbers sufficiently small, we can adjust the annular domain $A_{R_0,R_1}$ so that within $A_{R_0,R_1}$, the Gauss map remains free of poles and zeros. Accordingly,
   $$  f(z)=\frac{c_1^2z^2+\delta_1^2z^{-2}+2\epsilon_1 \delta_1 z^{-1}+2c_1\epsilon_1z}{z}
    $$
We compute $\psi_3 =z\phi_3$, where $\phi_3=fg$,
$$
\psi_3 (z) = (c_1z+\epsilon_1 + \delta_1 z^{-1}) (c_2 z^{-1}+\epsilon_2 +\delta_2 z)
$$
Given $g(z)$ set
    $$
    G_-(z)= c_1z+\epsilon_1 + \delta_1 z^{-1},\ \ \  G_+(z)=c_2 z^{-1}+\epsilon_2 +\delta_2 z 
    $$
    The computation yields:
    $$
    F_-(z)=G_-^2(z)=c_1^2z^2+\delta_1^2z^{-2}+2\epsilon_1 \delta_1 z^{-1}+2c_1\epsilon_1z,
    $$
   
    $$
    F_+(z)=G_+^2(z)=c_2^2z^{-2}+\delta_2^2z^2+2\epsilon_2 \delta_2 z+2c_2\epsilon_2z^{-1}
    $$
    Note that the mappings $F_-$ and $F_+$ satisfy
    $$
    \int_0^{2\pi} F_- d\theta =0,\ \int_0^{2\pi} F_+ d\theta =0
    $$
    This ensures that the new data results in a well-defined map on the adjusted domain $A_{R_0,R_1}$, and it also implies that the new surface maintains a vertical flux.
   \medskip
   
    The mappings $g(z)$ and $f(z)$ provide the parameters for the following Weierstrass representation:
    $$
    \bold{X}(p) = Re \int_{z,z_0} (\phi_1,\phi_2,\phi_3)dz
    $$
    where,
    $$
    \begin{cases}
        \phi_1=\frac{f}{2}(1-g^2) \\
        \phi_2=\frac{if}{2} (1+g^2) \\
        \phi_3 =fg
    \end{cases}
    $$
    Recall that the mapping $\bold{X}(p)$ is well-defined on $A_{R_0,R_1}$. 
    \medskip
    
    Let $\widetilde{\Sigma} = \bold{X} (A_{R_0,R_1})  \subset \mathbb{R}^3$
    represent the minimal surface in $\mathbb{R}^3$ parameterized by $\bold{X}$. Now, let $\Omega'$ be an open slab delimited by two parallel planes $P'_{\pm}$, where $P'_{\pm}=\{ x_3=h'_{\pm} \}$. Choose $0<h'_+ < h_+$ as a small number and $h'_-=-h'_+$, consequently, $\Omega' \subset \Omega$. For $h'_+$  sufficiently small, the associated Gauss map of $\bold{X}$ possesses a well-defined winding number, indicative of the number of times the circles foliating the surface $\widetilde{\Sigma}$ wind around it.
    We can confirm that this winding number remains two by directly evaluating the integral $\int_0^{2\pi} \frac{\partial}{\partial \theta} arg F(re^{i\theta}) d\theta$. Here, $\epsilon_j$ and $\delta_j$ are chosen to be sufficiently small, ensuring that the integral can be approximated by $\int_0^{2\pi} \frac{\partial}{\partial \theta} arg (c_j^2r^2e^{2i\theta}) d\theta$, for $j=1,2$. \\
  Define
$$
\Sigma = \widetilde{\Sigma} \cap \Omega'
$$
Subsequently,
$$
\bold{X}^{-1}(\Sigma) = A
$$
for some topological annulus $A$. By the Riemann mapping theorem there exists a conformal diffeomorphism $T$ from some planar annulus $A_{R_2,R_3}$ onto $A$, for certain positive values $R_2<R_3$, and $A_{R_2,R_3} \subset A_{R_0,R_1}$. Consider the map $\bold{X}\circ T: A_{R_2,R_3}\rightarrow \mathcal{A}_V(\Omega',2)$, and note that for a sufficiently thin slab $\Omega'$, and subsequently, a thin annulus $A$, the map $T$ is close to identity. Hence, $\bold{X}\circ T :=\hat{\bold{X}}$ has the Weierstrass data that approximates that of $\bold{X}$, thus $\hat{\bold{X}}=\bold{X}\circ T(A_{R_2,R_3})\in \mathcal{A}_V(\Omega',2)$.  
Observing that, 
$$
\int_0^{2\pi} \hat{G}_-(z)d\theta =2\pi \hat{\epsilon}_1 \neq 0,
$$
$$
\int_0^{2\pi} \hat{G}_+(z)d\theta =2\pi \hat{\epsilon}_2 \neq 0,
$$
we can conclude from Lemma \ref{lem: A* implies int G =0} that $\hat{\bold{X}} \in \mathcal{A}_V(\Omega',2) \backslash \mathcal{A}^*(\Omega',2)$.\\
Finally, we verify that
$$
\begin{cases}
    \int_0^{2\pi}r^2\Delta \vert \hat{G}_-\vert^2 d\theta =4\int_0^{2\pi} \vert \hat{G}_-\vert^2 d\theta -8\pi \vert \hat{\epsilon}_1 \vert^2 \\
    \int_0^{2\pi}r^2\Delta \vert \hat{G}_+\vert^2 d\theta =4\int_0^{2\pi} \vert \hat{G}_+\vert^2 d\theta -8\pi \vert \hat{\epsilon}_2 \vert^2
\end{cases}
$$
 Consequently,
 $$
 L''(t)=4L(t)-8\pi \vert \hat{\epsilon} \vert^2
 $$
 Given that $\epsilon$ is non-zero, it follows that 
 $$ L''(t)<4L(t),\ \forall t.$$
\end{proof}

\end{document}
